\section{Aufbau des \"Ubertragers}
\label{chap:uebertrager}
In dieser Aufgabe wird die bestehende Schaltung um einen gleichspannungssperrenden Ausgangskondensator und den \"Ubertrager erweitert. Der Kondensator blockiert den Gleichspannungsanteil der MOSFET-Endstufe, sodass am \"Ubertrager nur der Wechselspannungsanteil anliegt. Der \"Ubertrager wird als Rahmenspule ausgef\"uhrt und auf eine Zielinduktivit\"at von $150\,\mu H$ ausgelegt.

\subsection{Schaltung und Berechnung}
Die Auslegung erfolgt mit der N\"aherungsformel f\"ur eine Rahmenspule:
\[
  L = N^2 \cdot d \cdot 10^{-6}
\]
Nach der Windungszahl $N$ umgestellt ergibt sich:
\[
  N = \sqrt{\frac{L}{d \cdot 10^{-6}}}
\]
Mit $L=150\,\mu H$ und $d=6{,}5\,cm$ folgt:
\[
  N = \sqrt{\frac{150\,\mu H}{0{,}065\,m \cdot 10^{-6}}} \approx 48 \,\text{Windungen}
\]
Anschlie\ss end wird die Schaltung gem\"a\ss{} \cref{fig:schaltung_endstufe} aus \cref{chap:Signalaufbereitung} aufgebaut.

\subsection{Messvorgang}
Der \"Ubertrager wird zun\"achst mit rund 48 Windungen gewickelt. Die Sekund\"arspule wird dabei in gleicher Weise ausgef\"uhrt. Danach wird die resultierende Induktivit\"at mit einem LCR-Meter gemessen und mit dem Sollwert verglichen.

\subsection{Messergebnis}
Bei der ersten Messung mit etwa 48 Windungen wurde eine Induktivit\"at von rund $280\,\mu H$ ermittelt. Nach einer Reduktion auf ca. 36 Windungen lag die gemessene Induktivit\"at bei etwa $147\,\mu H$.

\subsection{Interpretation}
Die theoretische Auslegung liefert zwar etwa 48 Windungen, in der praktischen Umsetzung f\"uhrte dieser Wert jedoch zu einer deutlich zu hohen Induktivit\"at. Durch die Anpassung auf rund 36 Windungen konnte mit ca. $147\,\mu H$ ein Wert nahe der geforderten $150\,\mu H$ erreicht werden. Der Serienkondensator erf\"ullt dabei seine Aufgabe, indem er den Gleichspannungsanteil sperrt und den \"Ubertrager ausschlie\ss lich mit Wechselspannung speist.
